\documentstyle[%


righttag%


,11pt%


,amssymb%


,russian%


,izvru%


]{amsart}





\newcommand{\tl}[3]{\medskip\noindent{\bf#1}\ #2 \dotfill #3 \par} 


\pagestyle{empty}


\begin{document}


%\tableofcontents


\par


\vskip 2cm


\par


\bigskip


\bigskip


\bigskip


\par


\centerline{\Large СОДЕРЖАНИЕ}


\bigskip


\bigskip


%%%%%%%%%% Use \newline %%%%%%%%%%%





\tl{Андреев К.В.}


{О внутренних геометриях многообразия плоских


образующих 6-мерной\newline квадрики}{3}


\tl{Билялов Р.Ф., Никитин Б.С.}


{Спиноры в произвольных реперах.


Ковариантная\newline производная и производная Ли спиноров}{9}


\tl{Вишневский В.В.}


{Многообразие с парой коммутативных


полукасательных структур}{19}


\tl{Волак Р.А.}{О слоях лагранжевых слоений}{27}


\tl{Лейко С.Г.}


{$P$-геодезические преобразования и их группы


в касательных расслоениях,\newline индуцированные 


конциркулярными преобразованиями базисного многообразия}{35}


\tl{Мирзоян В.А.}


{Об одном классе подмногообразий с параллельной


фундаментальной\newline формой высшего порядка}{46}


\tl{Прванович Милева}


{О вполне омбилических подмногообразиях,


погруженных в слабо\newline симметрическое риманово многообразие}{54}


\tl{Розенфельд Б.А.}


{Теоретико-групповой смысл пространств Нейфельда}{65}


\tl{Семенов К.В.} 


{О геометрии эволюционного уравнения третьего порядка}{67}


\tl{Шандра И.Г.}


{О геометрии касательного расслоения


над многообразием с псевдосвязностью\newline


и антикватернионных $f$-структурах}{75}


\tl{Шибяк А.}


{Деривационные уравнения грассманианов второго порядка}{87}





\bigskip


{\centerline {Краткие сообщения}}


\bigskip





\tl{Чешкова М.А.} 


{К геометрии центральной проекции


$\size{14}{18pt}\selectfont n$-поверхности


в евклидовом\newline пространстве $E^{n+m}$}{96}














\vfill


\noindent\raisebox{2pt}{\copyright} Казанский госудаpственный унивеpситет


\end{document}


